\documentclass[10pt,a4paper]{amsart}
\usepackage[margin=19mm]{geometry}
\usepackage{amsmath,amssymb,mathtools}
\usepackage[scr=rsfs,cal=euler]{mathalfa}
\usepackage{xfrac}
\usepackage[unicode,hidelinks,bookmarks=false]{hyperref}
\newcommand{\ie}{i.e.\ }
\newcommand{\cf}{cf.\ }
\newcommand{\resp}{respectively\ }
\newcommand{\Subsection}{\subsection}
\providecommand{\nobreakdash}{\nobreak\hbox{-}}
\setlength{\emergencystretch}{2em}
\multlinegap0pt
\usepackage{bm}
\usepackage{bbm}
\usepackage{dsfont}
\usepackage{xspace}
\usepackage{cancel}
\usepackage{mathtools}
\usepackage{amsthm}
\makeatletter
\@ifundefined{theorem}{\newtheorem{theorem}{Theorem}}{}
\@ifundefined{proposition}{\newtheorem{proposition}[theorem]{Proposition}}{}
\makeatother
\makeatletter
\newcommand{\C}{{\mathbb C}}
\newcommand{\PGL}{{\mathrm{PGL}}}
\newcommand{\R}{{\mathbb R}}
\expandafter\ifx\csname assumption\endcsname\relax
\newenvironment{assumption}{\begin{assumption0}\rm}{\end{assumption0}}
\fi
\newcommand{\cB}{{\mathcal B}}
\expandafter\ifx\csname claim\endcsname\relax
\newenvironment{claim}{\begin{claim0}}{\end{claim0}}
\fi
\expandafter\ifx\csname conjecture\endcsname\relax
\newenvironment{conjecture}{\begin{conjecture0}}{\end{conjecture0}}
\fi
\expandafter\ifx\csname corollary\endcsname\relax
\newenvironment{corollary}{\begin{corollary0}}{\end{corollary0}}
\fi
\expandafter\ifx\csname definition\endcsname\relax
\newenvironment{definition}{\begin{defn0}\rm}{\end{defn0}}
\fi
\expandafter\ifx\csname example\endcsname\relax
\newenvironment{example}{\begin{example0}\rm}{\end{example0}}
\fi
\expandafter\ifx\csname lemma\endcsname\relax
\newenvironment{lemma}{\begin{lemma0}}{\end{lemma0}}
\fi
\expandafter\ifx\csname notation\endcsname\relax
\newenvironment{notation}{\begin{notation0}}{\end{notation0}}
\fi
\expandafter\ifx\csname proposition\endcsname\relax
\newenvironment{proposition}{\begin{prop0}}{\end{prop0}}
\fi
\expandafter\ifx\csname question\endcsname\relax
\newenvironment{question}{\begin{question0}\rm}{\end{question0}}
\fi
\expandafter\ifx\csname remark\endcsname\relax
\newenvironment{remark}{\begin{remark0}\rm}{\end{remark0}}
\fi
\expandafter\ifx\csname theorem\endcsname\relax
\newenvironment{theorem}{\begin{thm0}}{\end{thm0}}
\fi
\makeatother
\title[Periods of modular forms and applications to the conjectures]{Periods of modular forms and applications to the conjectures of Oda and of Prasanna-Venkatesh}
\author{Xavier Guitart, Santiago Molina}
\date{Source: arXiv:2507.05021v2 (CC BY 4.0)}
\begin{document}
\maketitle

\noindent\emph{Source and licence.} Periods of modular forms and applications to the conjectures of Oda and of Prasanna-Venkatesh. Version arXiv:2507.05021v2 (CC BY 4.0), distributed under the Creative Commons Attribution 4.0 International licence, as recorded in the arXiv licence metadata for this exact version and shown on the abstract page https://arxiv.org/abs/2507.05021v2. Full rights notice: licenses/arxiv-2507.05021-CC-BY-4.0.txt.\par\smallskip

\noindent\textit{Editorial scope.} Counted unit: the Proposition whose source label is \texttt{acthatHforvarphi} in the article (source0.tex lines 1271--1281, source label \texttt{acthatHforvarphi}), delivered together with the article's own complete original proof of it (source0.tex lines 1282--1339, one \texttt{proof} environment of 58 lines). No mathematics is written, repaired or re-derived: statement and proof are the article's own text line for line. Required context retained uncounted: eqstar2 (lines 1167--1169). Every definition, display and label of those lines is retained unchanged. Omitted from this package and not redistributed: 1--1166, 1170--1270, 1340--2917. Nothing inside a retained excerpt is omitted or altered: every retained body is delivered exactly as the article's lines are written, so no omission receipt of any kind accompanies this package. Citations: the retained excerpts carry no citation command at all, so no bibliography entry of the article is transcribed and no reference is redistributed. Reference adaptation: none was needed and none was made; every reference inside the delivered unit resolves inside it, so no unresolved reference or citation is produced. Printed slips kept literal: no printed word, symbol, hypothesis or number of the counted statement or of its proof was altered. Layout and class: the article's own class is \texttt{amsart} with packages including inputenc, color, hyperref, cleveref, xy, amssymb, amsmath, geometry, stackrel, csquotes, babel, enumerate, pdfpages, graphicx; the wrapper is the lane's pinned \texttt{amsart} editorial wrapper at 10pt on A4, which restates the theorem-like environments the retained text needs and adds the article's own macro definitions that the retained text needs (\texttt{\textbackslash C}, \texttt{\textbackslash PGL}, \texttt{\textbackslash R}, \texttt{\textbackslash cB}), with extra wrapper packages (bm, bbm, dsfont, xspace, cancel, mathtools). The delivered numbering is the unit's own. Assets: no retained span contains a figure environment, a graphics-inclusion command, a PDF-page inclusion, a table or a TikZ picture; no image file is copied into this package, the package declares no asset dependency and no image file is redistributed with it. Licence: version arXiv:2507.05021v2 of this article is distributed under the Creative Commons Attribution 4.0 International licence, as recorded in the arXiv licence metadata for this exact version and shown on the abstract page https://arxiv.org/abs/2507.05021v2. A full rights notice is delivered at licenses/arxiv-2507.05021-CC-BY-4.0.txt. Characters: every retained excerpt is pure ASCII, so the delivered engine, XeLaTeX with its default Latin Modern text font, reports no missing character.
\begin{equation}\label{eqstar2}
g=u\left(\begin{array}{cc}r&x\\&r^{-1}\end{array}\right)\kappa(\alpha,\beta),\qquad \text{with } u\in\C^\times\;,x\in\C\;,r\in\R^\times,    
\end{equation}

\begin{proposition}\label{acthatHforvarphi}
    If $\chi(re^{i\theta})=r^{N_{\chi}}e^{i\lambda_\chi \theta}$ then the action of $\hat H$ on $\tilde \cB(\chi)$ is given by
\[
\hat H \varphi_n(\mu)=\lambda_{\chi}(N_{\chi}-1)\varphi_n(\mu_0)-(N_{\chi}+n)\varphi_{n+1}(\mu_1)+(n+\lambda_{\chi})(n-\lambda_{\chi})(n-N_{\chi}+1)\varphi_{n-1}(\mu_{-1}),
\]
where $\mu_0\in V(2n)$, $\mu_1\in V(2n+2)$ and $\mu_{-1}\in V(2n-2)$ are
\[
\mu_0(P):=\frac{1}{n(n+1)}\mu\left(
y\frac{\partial P}{\partial y}-x\frac{\partial P}{\partial x}\right),\quad\mu_1(P):=\frac{2}{(n+1)(2n+1)}\mu\left(\frac{\partial P}{\partial x\partial y}\right),\quad \mu_{-1}(P):=\frac{2}{n(2n+1)}\mu\left(xyP\right).
\] 
\end{proposition}

\begin{proof}
        If we consider $f:\PGL_2(\C)\rightarrow\C$ as a function with variables $s,r,\alpha,\beta$ by means of \eqref{eqstar2}, then
    \begin{eqnarray*}
    \hat H f(s,r,\alpha,\beta)&=&\frac{d}{dt}f(g\exp(t\hat H))\mid_{t=0}=\frac{d}{dt}f\left(\left(\begin{array}{cc}r&s\\&r^{-1}\end{array}\right)\left(\begin{array}{cc}\alpha&\beta\\-\bar \beta&\bar \alpha\end{array}\right)\left(\begin{array}{cc}
       e^{t}  & 0 \\
        0 & e^{-t}
    \end{array}\right)\right)\mid_{t=0}\\
    &=&\frac{d}{dt}f\left(\left(\begin{array}{cc}r&s\\&r^{-1}\end{array}\right)\left(\begin{array}{cc}e^t\alpha&e^{-t}\beta\\-e^t\bar \beta&e^{-t}\bar \alpha\end{array}\right)\right)\mid_{t=0}=\frac{d}{dt}f\left(\left(\begin{array}{cc}rR^{-1}&sR+rA\\&r^{-1}R\end{array}\right)\left(\begin{array}{cc}\frac{e^{-t}\alpha}{R}&\frac{e^{t}\beta}{R}\\\frac{-e^t\bar \beta}{R}&\frac{e^{-t}\bar \alpha}{R}\end{array}\right)\right)\mid_{t=0}\\
    &=&\frac{d}{dt}f(sR+rA,rR^{-1},e^{-t}\alpha R^{-1},e^{t}\beta R^{-1})\mid_{t=0}
    \end{eqnarray*}
    where $R^2=R(t)^2=e^{2t}|\beta|^2+e^{-2t}|\alpha|^2$ and $A=A(t)=\alpha\beta R^{-1}(e^{-2t}-e^{2t})$.
    Since $R(0)=1$, $R'(0)=|\beta|^2-|\alpha|^2$, $A(0)=0$ and $A'(0)=-4\alpha\beta$, we conclude
    \begin{equation}\label{acthatH}
        \hat H f=(s(|\beta|^2-|\alpha|^2)-4r\alpha\beta)\frac{\partial f}{\partial s}+(\bar s(|\beta|^2-|\alpha|^2)-4r\bar\alpha\bar\beta)\frac{\partial f}{\partial \bar s}-r(|\beta|^2-|\alpha|^2)\frac{\partial f}{\partial r}-2\alpha|\beta|^2\frac{\partial f}{\partial \alpha}-2\bar \alpha|\beta|^2\frac{\partial f}{\partial \bar\alpha}+2\beta|\alpha|^2\frac{\partial f}{\partial \beta}+2\bar\beta|\alpha|^2\frac{\partial f}{\partial \bar\beta}.
    \end{equation}

If we write $\chi(re^{i\theta})=r^Ne^{i\lambda\theta}$,
then we have by definition
\begin{equation}\label{generalvarphin}
\varphi_n(\mu)\left(\left(\begin{array}{cc}r&s\\&r^{-1}\end{array}\right)\kappa(\alpha,\beta)\right):=r^{2N}\cdot\mu\left(
P_{n+\lambda,n-\lambda}\right);\qquad P_{a,b}:=\left|\begin{array}{cc}\alpha&\beta\\x&y\end{array}\right|^{a}\left|\begin{array}{cc}-\bar\beta&\bar\alpha\\x&y\end{array}\right|^{b}.
\end{equation}
Hence, we obtain by \eqref{acthatH}
\begin{eqnarray*}
\hat H\varphi_n(\mu)(\alpha,\beta)&=&\mu\left(2N(|\alpha|^2-|\beta|^2)P_{n+\lambda,n-\lambda}-2(n+\lambda)\alpha|\beta|^2yP_{n-1+\lambda,n-\lambda}+2(n-\lambda)\bar\alpha|\beta|^2xP_{n+\lambda,n-1-\lambda}-\right.\\
&&\left.-2(n+\lambda)\beta|\alpha|^2xP_{n-1+\lambda,n-\lambda}-2(n-\lambda)\bar\beta|\alpha|^2yP_{n+\lambda,n-1-\lambda}\right)\\
&=&\mu\left(2N(|\alpha|^2-|\beta|^2)P_{n+\lambda,n-\lambda}+2(n+\lambda)(-\alpha|\beta|^2y-\beta|\alpha|^2x)P_{n-1+\lambda,n-\lambda}+\right.\\
&&\left.+2(n-\lambda)(\bar\alpha|\beta|^2x-\bar\beta|\alpha|^2y)P_{n+\lambda,n-1-\lambda}\right)\\
&=&\mu\left(2N(|\alpha|^2-|\beta|^2)P_{n+\lambda,n-\lambda}+(n+\lambda)(-\alpha y-\beta x+(\beta x-\alpha y)(|\beta|^2-|\alpha|^2))P_{n-1+\lambda,n-\lambda}+\right.\\
&&\left.+(n-\lambda)(\bar\alpha x-\bar\beta y+(\bar\alpha x+\bar\beta y)(|\beta|^2-|\alpha|^2))P_{n+\lambda,n-1-\lambda}\right)\\
&=&\mu\left((2N+2n)(|\alpha|^2-|\beta|^2)P_{n+\lambda,n-\lambda}+(n+\lambda)(-\alpha y-\beta x)P_{n-1+\lambda,n-\lambda}+\right.\\
&&\left.+(n-\lambda)(\bar\alpha x-\bar\beta y)P_{n+\lambda,n-1-\lambda}\right)\\
&=&\mu\left((2N+2n)(|\alpha|^2-|\beta|^2)P_{n+\lambda,n-\lambda}+(n+\lambda)2xyP_{n-1+\lambda,n-1-\lambda}+2\lambda(\bar\beta y-\bar\alpha x) P_{n+\lambda,n-1-\lambda}\right),
\end{eqnarray*}
where the last equality follows from $(-\beta x-\alpha y)P_{a-1,b}=2xyP_{a-1,b-1}+(\bar\beta y-\bar\alpha x) P_{a,b-1}$ deduced from the relation
\begin{equation}\label{keyequationPab}
    -\beta P_{a,b+1}=yP_{a,b}-\bar\alpha P_{a+1,b};\qquad -\alpha P_{a,b+1}=xP_{a,b}+\bar\beta P_{a+1,b}.
\end{equation}
We compute similarly
\begin{eqnarray*}
y\frac{\partial P_{a,b}}{\partial y}-x\frac{\partial P_{a,b}}{\partial x}&=&y(a\alpha P_{a-1,b}-b\bar\beta P_{a,b-1})-x(-a\beta P_{a-1,b}-b\bar\alpha P_{a,b-1})=a(y\alpha+x\beta)P_{a-1,b}+b(x\bar\alpha -y\bar\beta )P_{a,b-1}\\
&=&-2axyP_{a-1,b-1}-(a+b)(\bar\beta y-\bar\alpha x) P_{a,b-1}.
\end{eqnarray*}
Moreover, using the relations \eqref{keyequationPab}
one obtains
\begin{eqnarray*}
\frac{\partial^2 P_{a,b}}{\partial y\partial x}&=&-a\beta ((a-1)\alpha P_{a-2,b}-b\bar\beta P_{a-1,b-1})-b\bar\alpha (a\alpha P_{a-1,b-1}-(b-1)\bar\beta P_{a,b-2})\\
&=&-a (a-1)\beta\alpha P_{a-2,b}+ab(|\beta|^2 -|\alpha|^2) P_{a-1,b-1}+b(b-1)\bar\alpha\bar\beta P_{a,b-2}\\
&=&a (a-1)(y\alpha P_{a-2,b-1}-|\alpha|^2P_{a-1,b-1})+ab(|\beta|^2 -|\alpha|^2) P_{a-1,b-1}+b(b-1)(-|\alpha|^2 P_{a-1,b-1}-x\bar\alpha P_{a-1,b-2})\\
&=&\frac{(a+b)(a+b-1)}{2}(|\beta|^2 -|\alpha|^2) P_{a-1,b-1}-a(a-1)yxP_{a-2,b-2}-\frac{a(a-1)-b(b-1)}{2}(y\bar\beta -x\bar\alpha) P_{a-1,b-2}
\end{eqnarray*}
Thus,
\begin{eqnarray*}
\hat H\varphi_n(\mu)(\alpha,\beta)&=&\mu\left(\frac{-2(N+n)}{(n+1)(2n+1)}\frac{\partial^2 P_{n+1+\lambda,n+1-\lambda}}{\partial y\partial x}+\frac{\lambda(N-1)}{n(n+1)}\left(y\frac{\partial P_{n+\lambda,n-\lambda}}{\partial y}-x\frac{\partial P_{n+\lambda,n-\lambda}}{\partial x}\right)+\right.\\
&&\left.+\frac{2(n+\lambda)(n-\lambda)(n-N+1)}{n(2n+1)}xyP_{n-1+\lambda,n-1-\lambda}\right),
\end{eqnarray*}
and the result follows.
\end{proof}

\end{document}
